\begin{theorem}[Three-point flatness]
Let $E$ be a finite-dimensional real inner-product space, write $B_\rho=B(0,\rho)$
for the open ball of radius $\rho$ about the origin, and let $L\in\R$. Let
$U_i:E\to\R$, $i\in\N$, be infinity-harmonic on $B_1$ in the sense of
\noderef{IsInfHarmonicOn} and satisfy the common Lipschitz bound
\[
  |U_i(x)-U_i(y)|\le L|x-y|\qquad(x,y\in B_1,\ i\in\N).
\]
Let $q_i\in E$ and $b_i\in\R$, and write $\ell_i(x)=\langle q_i,x\rangle+b_i$ for
the corresponding affine functions. Assume that $U_i-\ell_i\to0$ uniformly on
$B_1$, that is,
\[
  \text{for every }\varepsilon>0\text{ there is }i_0\text{ with }
  |U_i(x)-\ell_i(x)|\le\varepsilon
  \quad(i\ge i_0,\ x\in B_1).
\]
Then the midpoint defect of $U_i$, divided by the endpoint separation, tends to
zero uniformly on $B_{1/4}$: for every $\varepsilon>0$ there is $i_0$ such that
\[
  \Bigl|U_i(x)+U_i(y)-2U_i\bigl(\tfrac12(x+y)\bigr)\Bigr|\le\varepsilon|x-y|
  \qquad\bigl(i\ge i_0,\ x,y\in B_{1/4},\ x\ne y\bigr).
\]
\end{theorem}
\begin{proof}
\emph{Step 0: the degenerate case and the normalizations.}
If every element of $E$ is $0$ then there are no points $x,y$ with $x\ne y$ and
the conclusion holds vacuously with $i_0=0$. So assume $E\ne\{0\}$. Then $B_1$
contains two distinct points $x\ne y$, and the Lipschitz hypothesis applied to
them gives $0\le|U_i(x)-U_i(y)|\le L|x-y|$ with $|x-y|>0$, so
\[
  L\ge0 .
\]

Write $w_i=U_i-\ell_i$. For $\varepsilon>0$ let $\mathcal E(\varepsilon)$ denote
an index $i_0$ supplied by the flatness hypothesis for that $\varepsilon$, so
that
\[
  |w_i(x)|\le\varepsilon\qquad\bigl(i\ge\mathcal E(\varepsilon),\ x\in B_1\bigr).
\]

\emph{The slopes $q_i$ are bounded.} Let $i\ge\mathcal E(1)$ and let $e\in E$
with $|e|=1$. The points $\pm e/2$ lie in $B_1$, and
$\ell_i(e/2)-\ell_i(-e/2)=\langle q_i,e\rangle$, so
\[
  |\langle q_i,e\rangle|
   \le\bigl|U_i(e/2)-U_i(-e/2)\bigr|+|w_i(e/2)|+|w_i(-e/2)|
   \le L|e|+2=L+2 .
\]
Taking $e=q_i/|q_i|$ when $q_i\ne0$ gives $|q_i|\le L+2$ (the bound being
trivial otherwise). Consequently, for $i\ge\mathcal E(1)$ and $x,y\in B_1$,
\[
  |w_i(x)-w_i(y)|\le|U_i(x)-U_i(y)|+|\langle q_i,x-y\rangle|
   \le\bigl(L+(L+2)\bigr)|x-y|,
\]
so with
\[
  C:=2L+2\ge2>0
\]
every $w_i$ with $i\ge\mathcal E(1)$ is $C$-Lipschitz on $B_1$, and
$|q_i|\le L+2\le C$.

\emph{Negation is admissible.} We record, for use below, the following scaling
property. Let $\Omega\subseteq E$ be open, let $u:E\to\R$ and let $t<0$. First,
for any $C^2$ function $\phi$ near a point $z$ and any $s\in\R$,
\begin{equation}\label{eq:tpf-scale}
  \Dinf(s\phi)(z)=s^3\,\Dinf\phi(z),
\end{equation}
because $\nabla(s\phi)(z)=s\nabla\phi(z)$ and $D^2(s\phi)(z)=sD^2\phi(z)$, so
that $D^2(s\phi)(z)[\nabla(s\phi)(z)]
=s\,D^2\phi(z)\bigl(s\nabla\phi(z),s\nabla\phi(z)\bigr)=s^3D^2\phi(z)[\nabla\phi(z)]$,
using the definition of $\Dinf$ in \noderef{InfLap} and the bilinearity of the
second derivative. Now suppose $u$ is an infinity supersolution on $\Omega$ in
the sense of \noderef{IsInfSupersolutionOn}, and let $\phi$ be $C^2$ near
$z\in\Omega$ with $tu-\phi$ having a local maximum on $\Omega$ at $z$. Dividing
the inequality $tu(x)-\phi(x)\le tu(z)-\phi(z)$ by $t<0$ reverses it and gives
$u(x)-\phi(x)/t\ge u(z)-\phi(z)/t$ for $x$ near $z$ in $\Omega$; that is,
$\phi/t$ is a $C^2$ lower test for $u$ at $z$, so $\Dinf(\phi/t)(z)\le0$ by
\noderef{IsInfSupersolutionOn}. By \eqref{eq:tpf-scale} with $s=1/t$ this reads
$t^{-3}\Dinf\phi(z)\le0$, and $t^{-3}<0$, so $\Dinf\phi(z)\ge0$. As $\Omega$ is
open and $tu$ is continuous on $\Omega$, this says that $tu$ is an infinity
subsolution on $\Omega$ in the sense of \noderef{IsInfSubsolutionOn}. The
symmetric argument, starting from $u$ an infinity subsolution and a lower test
for $tu$, shows that $tu$ is then an infinity supersolution. In particular, if
$u$ is infinity-harmonic on $\Omega$ then so is $-u$, and if $u$ is
infinity-harmonic on $\Omega$ then $-2u$ is an infinity subsolution on $\Omega$.

\emph{Step 1: the contradiction hypothesis.}
Suppose the conclusion fails. Then there is $\delta>0$ such that for every
$i_0$ there are $i\ge i_0$ and $x,y\in B_{1/4}$ with $x\ne y$ and
\[
  \Bigl|U_i(x)+U_i(y)-2U_i\bigl(\tfrac12(x+y)\bigr)\Bigr|>\delta|x-y| .
\]
Since $\ell_i$ is affine and $\tfrac12(x+y)$ is the midpoint,
$\ell_i(x)+\ell_i(y)-2\ell_i(\tfrac12(x+y))=0$, so the displayed quantity is
unchanged when $U_i$ is replaced by $w_i$. Choosing $i_0=1,2,\dots$ in turn and
discarding repetitions produces a strictly increasing sequence of indices
$i_1<i_2<\cdots$ and points $x_k,y_k\in B_{1/4}$, $x_k\ne y_k$, with
\begin{equation}\label{eq:tpf-violate}
  \bigl|w_{i_k}(x_k)+w_{i_k}(y_k)-2w_{i_k}(c_k)\bigr|>\delta|x_k-y_k|,
  \qquad c_k=\tfrac12(x_k+y_k).
\end{equation}
Passing to a further subsequence we may assume that the sign of
$w_{i_k}(x_k)+w_{i_k}(y_k)-2w_{i_k}(c_k)$ is the same for every $k$. If that
sign is negative, replace the whole family $(U_i,q_i,b_i)$ by
$(-U_i,-q_i,-b_i)$: by the scaling property of Step 0 each $-U_i$ is again
infinity-harmonic on $B_1$, the Lipschitz bound with the same $L$ is unchanged,
the flatness hypothesis is unchanged, $w_i$ is replaced by $-w_i$, and the
quantity in \eqref{eq:tpf-violate} changes sign. We may therefore assume
\begin{equation}\label{eq:tpf-positive}
  w_{i_k}(x_k)+w_{i_k}(y_k)-2w_{i_k}(c_k)>\delta|x_k-y_k|\qquad(k\ge1).
\end{equation}
Note also that, by the $C$-Lipschitz bound and $|x_k-c_k|=|y_k-c_k|=\tfrac12
|x_k-y_k|$, the left side of \eqref{eq:tpf-positive} is at most
\[
  \bigl(w_{i_k}(x_k)-w_{i_k}(c_k)\bigr)+\bigl(w_{i_k}(y_k)-w_{i_k}(c_k)\bigr)
   \le C\cdot\tfrac12|x_k-y_k|+C\cdot\tfrac12|x_k-y_k|=C|x_k-y_k|
\]
whenever $i_k\ge\mathcal E(1)$. Choosing one such $k$ --- possible because
$i_k\to\infty$ --- and dividing by $|x_k-y_k|>0$ gives
\begin{equation}\label{eq:tpf-deltaC}
  \delta\le C .
\end{equation}

\emph{Step 2: the two auxiliary profiles.}
For $0<r\le3/2$ set
\[
  f(r)=r-\tfrac18r^{3/2},\qquad
  f'(r)=1-\tfrac3{16}\sqrt r,\qquad
  f''(r)=-\frac3{32\sqrt r},
\]
so $f$ is $C^2$ on $(0,\infty)$. Since $\sqrt r\le\sqrt{3/2}$ on $(0,3/2]$,
\begin{equation}\label{eq:tpf-f}
  \nu r\le f(r)\le r,\qquad \nu:=1-\tfrac18\sqrt{3/2}>0,
  \qquad 0<f'(r)\le1 .
\end{equation}

Next we build the localizing factor. Let
\[
  P(\tau)=\begin{cases}0,&\tau\le0,\\
    6\tau^5-15\tau^4+10\tau^3,&0\le\tau\le1,\\ 1,&\tau\ge1,\end{cases}
\]
a $C^2$ function on $\R$: on $[0,1]$ one computes $P(0)=0$, $P(1)=1$,
$P'(\tau)=30\tau^2(\tau-1)^2$ and $P''(\tau)=60\tau(2\tau-1)(\tau-1)$, so
$P,P',P''$ all match the constant branches in value at $\tau=0$ and $\tau=1$;
moreover $P'\ge0$, so $0\le P\le1$. Fix
\[
  B_{\mathrm{out}}:=\frac{C+\tfrac32C^2}{\nu}+\frac{\delta}{2}+1
\]
and define $B:E\to\R$ by
\[
  B(c)=\frac\delta2+\Bigl(B_{\mathrm{out}}-\frac\delta2\Bigr)
        P\Bigl(\frac{16|c|^2-1}{3}\Bigr).
\]
Since $c\mapsto|c|^2$ is a continuous quadratic form, hence $C^\infty$, and $P$
is $C^2$, the function $B$ is $C^2$ on $E$. Moreover
$B_{\mathrm{out}}-\delta/2>0$ and $0\le P\le1$ give
\begin{equation}\label{eq:tpf-B}
  \frac\delta2\le B\le B_{\mathrm{out}},\qquad
  B\equiv\frac\delta2\ \text{on }\overline B(0,1/4),\qquad
  B\equiv B_{\mathrm{out}}\ \text{where }|c|\ge\tfrac12,
\end{equation}
because $|c|\le1/4$ makes the argument of $P$ at most $0$ and $|c|\ge1/2$ makes
it at least $1$. Finally put
\[
  K_1=\max_{\overline B(0,3/4)}|DB|,\qquad
  K_2=\max_{\overline B(0,3/4)}\|D^2B\| ,
\]
finite because $DB$ and $D^2B$ are continuous on the compact set
$\overline B(0,3/4)$ (preamble).

\emph{Step 3: the maximization.}
Fix $k$, write $i=i_k$ and assume $i\ge\mathcal E(1)$. For a triple
$(x,y,z)\in\bigl(\overline B(0,3/4)\bigr)^3$ with $x\ne y$ put
\[
  c=\tfrac12(x+y),\qquad h=z-c,\qquad r=|x-y|\in(0,3/2],
\]
and
\[
  \Phi_i(x,y,z)=w_i(x)+w_i(y)-2w_i(z)-B(c)f(r)-\frac{|h|^2}{r^2} .
\]
Note $c\in\overline B(0,3/4)$ by convexity, and all three points lie in $B_1$.

\emph{(3a) An upper bound.} Writing
$w_i(x)+w_i(y)-2w_i(z)=\bigl(w_i(x)-w_i(c)\bigr)+\bigl(w_i(y)-w_i(c)\bigr)
+2\bigl(w_i(c)-w_i(z)\bigr)$ and using the $C$-Lipschitz bound together with
$|x-c|=|y-c|=r/2$ and $|c-z|=|h|$,
\[
  w_i(x)+w_i(y)-2w_i(z)\le Cr+2C|h| .
\]
Since $2C|h|-|h|^2/r^2\le C^2r^2$ (the quadratic $s\mapsto2Cs-s^2/r^2$ is
maximal at $s=Cr^2$ with value $C^2r^2$), we get
\begin{equation}\label{eq:tpf-upper}
  \Phi_i(x,y,z)\le Cr+C^2r^2-B(c)f(r)
   \le Cr+C^2r^2=:rC',\qquad C':=C+\tfrac32C^2,
\end{equation}
using $B\ge0$, $f\ge0$ and $r\le3/2$ in the last step.

\emph{(3b) The supremum is positive.} Take $x=x_k$, $y=y_k$, $z=c_k$, so that
$h=0$ and $r=r_k:=|x_k-y_k|>0$. Then $c=c_k\in\overline B(0,1/4)$, since
$x_k,y_k\in B_{1/4}$, so $B(c_k)=\delta/2$ by \eqref{eq:tpf-B}, and by
\eqref{eq:tpf-positive} and $f(r)\le r$,
\[
  \Phi_i(x_k,y_k,c_k)>\delta r_k-\frac\delta2 f(r_k)
   \ge\delta r_k-\frac\delta2 r_k=\frac{\delta r_k}{2}=:\beta_k>0 .
\]
Hence
\[
  M_i:=\sup\bigl\{\Phi_i(x,y,z):x,y,z\in\overline B(0,3/4),\ x\ne y\bigr\}
\]
satisfies $M_i\ge\beta_k>0$, and $M_i\le\tfrac32C'<\infty$ by
\eqref{eq:tpf-upper}.

\emph{(3c) The supremum is attained, at a triple with $r$ bounded below.} Put
\[
  \mathcal K=\Bigl\{(x,y,z)\in\bigl(\overline B(0,3/4)\bigr)^3:
    |x-y|\ge\frac{\beta_k}{2C'}\Bigr\},
\]
a closed and bounded, hence compact, subset of $E^3$ (preamble). Any admissible
triple with $r<\beta_k/(2C')$ has $\Phi_i\le rC'<\beta_k/2<\beta_k\le M_i$ by
\eqref{eq:tpf-upper}, so such triples do not affect the supremum: $M_i$ equals
the supremum of $\Phi_i$ over $\mathcal K$, which is nonempty because
$\beta_k/(2C')=\delta r_k/(4C')\le r_k$ by \eqref{eq:tpf-deltaC} and
$C'\ge C\ge\delta$, so $(x_k,y_k,c_k)\in\mathcal K$. On $\mathcal K$ we have
$|x-y|>0$, so $\Phi_i$ is continuous there, and a continuous function on a
nonempty compact set attains its maximum (preamble). Let $(x,y,z)$ be a
maximizing triple, so
\begin{equation}\label{eq:tpf-max}
  \Phi_i(x,y,z)=M_i\ge\beta_k>0 .
\end{equation}
By \eqref{eq:tpf-upper} and \eqref{eq:tpf-max}, $rC'\ge\beta_k$, i.e.
$r\ge\beta_k/C'>\beta_k/(2C')$: the maximizing triple is strictly inside the
constraint defining $\mathcal K$.

\emph{Step 4: the geometry at the maximizer.}
Keep the maximizing triple $(x,y,z)$ of Step 3 and its $c,h,r$.

\emph{(4a) $r$ is small when $i$ is large.} Let $\varepsilon>0$ and suppose
$i\ge\mathcal E(\varepsilon)$, so $|w_i|\le\varepsilon$ on $B_1$. From
\eqref{eq:tpf-max}, $B\ge\delta/2$ and $f(r)\ge\nu r$,
\[
  0<\Phi_i(x,y,z)\le4\varepsilon-B(c)f(r)\le4\varepsilon-\frac{\delta\nu r}{2},
\]
whence
\begin{equation}\label{eq:tpf-rsmall}
  r<\frac{8\varepsilon}{\delta\nu} .
\end{equation}

\emph{(4b) $|h|\le2Cr^2$.} The triple $(x,y,c)$ is admissible and lies in
$\mathcal K$ (it has the same $x,y$), so $\Phi_i(x,y,z)\ge\Phi_i(x,y,c)$. The
terms $w_i(x)+w_i(y)$ and $B(c)f(r)$ are identical in the two, and the penalty
vanishes at $z=c$, so
\[
  -2w_i(z)-\frac{|h|^2}{r^2}\ge-2w_i(c),
  \qquad\text{i.e.}\qquad
  \frac{|h|^2}{r^2}\le2\bigl(w_i(c)-w_i(z)\bigr)\le2C|h| .
\]
Hence $|h|^2\le2C|h|r^2$, and therefore
\begin{equation}\label{eq:tpf-h}
  |h|\le2Cr^2
\end{equation}
(trivially if $h=0$, and after dividing by $|h|>0$ otherwise).

\emph{(4c) $|c|<1/2$.} If $|c|\ge1/2$ then $B(c)=B_{\mathrm{out}}$ by
\eqref{eq:tpf-B}, so \eqref{eq:tpf-upper}, $f(r)\ge\nu r$ and $r\le3/2$ give
\[
  \Phi_i(x,y,z)\le Cr+C^2r^2-B_{\mathrm{out}}\nu r
   \le r\Bigl(C+\tfrac32C^2-B_{\mathrm{out}}\nu\Bigr)<0
\]
by the choice of $B_{\mathrm{out}}$, contradicting \eqref{eq:tpf-max}. So
$|c|<1/2$.

\emph{(4d) The maximizer is interior.} Assume in addition
\begin{equation}\label{eq:tpf-r1}
  r<\tfrac12\qquad\text{and}\qquad 2Cr^2<\tfrac14 .
\end{equation}
Then $|x|\le|c|+\tfrac12r<\tfrac12+\tfrac14=\tfrac34$ and likewise
$|y|<3/4$, while $|z|\le|c|+|h|<\tfrac12+\tfrac14=\tfrac34$ by
\eqref{eq:tpf-h}. So $(x,y,z)$ lies in the open set $\bigl(B_{3/4}\bigr)^3$.
Together with $r>\beta_k/(2C')$ and the continuity of $(x',y',z')\mapsto|x'-y'|$,
there is a neighbourhood $\mathcal N$ of $(x,y,z)$ in $E^3$ contained in
$\bigl(B_{3/4}\bigr)^3\cap\{|x'-y'|>\beta_k/(2C')\}\subseteq\mathcal K$. Hence
\begin{equation}\label{eq:tpf-localmax}
  \Phi_i\le\Phi_i(x,y,z)\quad\text{on }\mathcal N:
  \qquad(x,y,z)\ \text{is a local maximum of}\ \Phi_i\ \text{on}\ E^3 .
\end{equation}

\emph{Step 5: the sum rule.}
Set $e=(x-y)/r$, a unit vector, $q=q_i$, and
\[
  \Psi_0(x',y',z')=\ell_i(x')+\ell_i(y')-2\ell_i(z')
    +B(c')f(r')+\frac{|h'|^2}{(r')^2},
\]
where $c'=\tfrac12(x'+y')$, $h'=z'-c'$, $r'=|x'-y'|$. This formula defines
$\Psi_0$ on the open set $\mathcal O=\{(x',y',z')\in E^3:x'\ne y'\}$ only, since
$f(r')$ and $|h'|^2/(r')^2$ are undefined at $r'=0$. On $\mathcal O$ the maps
$(x',y')\mapsto c'$ and
$(x',y',z')\mapsto h'$ are linear, $r'$ is $C^\infty$ (being the norm away from
$0$), $f$ is $C^2$ on $(0,\infty)$ and $B$ is $C^2$; so $\Psi_0$ is $C^2$ on
$\mathcal O$, and $(x,y,z)\in\mathcal O$ because $r>\beta_k/(2C')>0$.

\emph{The total test.} \noderef{SumRuleJets} takes a test defined on all of
$V=\ell^2(E\oplus E\oplus E)$, so $\Psi_0$ must be extended; we name one
extension rather than leave it implicit. Put $\rho=\beta_k/(4C')>0$, so that
$r>\beta_k/(2C')=2\rho$ and, by (4d),
\begin{equation}\label{eq:tpf-cutoff-nbhd}
  \mathcal N\subseteq\bigl\{(x',y',z'):|x'-y'|>2\rho\bigr\}\subseteq\mathcal O .
\end{equation}
Define $\Psi:V\to\R$ by truncation outside that slab:
\[
  \Psi(x',y',z')=\begin{cases}
    \Psi_0(x',y',z')&|x'-y'|>2\rho,\\
    0&|x'-y'|\le2\rho .
  \end{cases}
\]
This is defined on all of $V$ --- the first branch is legitimate because
$|x'-y'|>2\rho>0$ forces $x'\ne y'$, so $(x',y',z')\in\mathcal O$ --- and
\begin{equation}\label{eq:tpf-psi-agrees}
  \Psi=\Psi_0\quad\text{on the open set}\quad
  \bigl\{(x',y',z'):|x'-y'|>2\rho\bigr\}\supseteq\mathcal N ,
\end{equation}
which contains $(x,y,z)$ because $r>2\rho$. Hence $\Psi$ is $C^2$ near
$(x,y,z)$, and $D\Psi(x,y,z)=D\Psi_0(x,y,z)$ and
$D^2\Psi(x,y,z)=D^2\Psi_0(x,y,z)$, two functions agreeing on a neighbourhood of
a point having the same derivatives there. So every derivative computation below
may be carried out on the formula for $\Psi_0$, and we write $\Psi$ for both
from here on.

Two remarks on why this suffices, and why nothing smoother is wanted. First,
\noderef{SumRuleJets} imposes no regularity on its test away from the point: its
two hypotheses involving $\Psi$ are that $\Psi$ be $C^2$ \emph{near}
$\widehat Z$ and that $\sum_jg_j-\Psi$ have a \emph{local} maximum at
$\widehat Z$, both local conditions, and its conclusion involves $\Psi$ only
through $D_j\Psi(\widehat Z)$ and $A=D^2\Psi(\widehat Z)$. The truncation above
is therefore enough even though it is discontinuous across
$\{|x'-y'|=2\rho\}$, and in particular no smooth cutoff and no existence
statement for one is needed. Second, for the same reason nothing depends on the
choice: any map $V\to\R$ agreeing with $\Psi_0$ on some neighbourhood of
$(x,y,z)$ would serve equally, and this one is named only so that no extension
has to be invented later.

With
\[
  g_1=g_2=U_i,\qquad g_3=-2U_i,
\]
each continuous on the open set $B_1$ (by \noderef{IsInfHarmonicOn}, which
includes continuity on $B_1$), we have, on $\mathcal N$,
\[
  g_1(x')+g_2(y')+g_3(z')-\Psi(x',y',z')=\Phi_i(x',y',z'),
\]
since $U_i-\ell_i=w_i$ and, by \eqref{eq:tpf-psi-agrees}, $\Psi=\Psi_0$ there;
the right-hand side is defined on $\mathcal N$ because
$\mathcal N\subseteq\bigl(B_{3/4}\bigr)^3$ and $x'\ne y'$ on $\mathcal N$ by
\eqref{eq:tpf-cutoff-nbhd}. By \eqref{eq:tpf-localmax} the right-hand side is
$\le\Phi_i(x,y,z)$ on $\mathcal N$, so $\sum_jg_j-\Psi$ has a local maximum at
$(x,y,z)$ on $V$, whose components lie in $B_{3/4}\subseteq B_1$.

Differentiating $\Psi$ at $(x,y,z)$ gives, with $\partial c'/\partial
x'=\partial c'/\partial y'=\tfrac12\mathrm{id}$, $\partial h'/\partial
z'=\mathrm{id}$, $\partial h'/\partial x'=\partial h'/\partial
y'=-\tfrac12\mathrm{id}$, $\partial r'/\partial x'=e$ and $\partial
r'/\partial y'=-e$,
\begin{equation}\label{eq:tpf-grads}
\begin{gathered}
  p_1:=D_{x}\Psi=q+\tfrac12f(r)DB(c)-\frac{h}{r^2}+ae,\qquad
  p_2:=D_{y}\Psi=q+\tfrac12f(r)DB(c)-\frac{h}{r^2}-ae,\\
  D_{z}\Psi=-2p_0\ \text{with}\ p_0:=q-\frac{h}{r^2},
  \qquad a:=B(c)f'(r)-\frac{2|h|^2}{r^3}.
\end{gathered}
\end{equation}
Indeed $D_z\Psi=-2q+2h/r^2=-2p_0$; and in the $x'$ variable the three pieces
contribute $q$, then
$\tfrac12f(r)DB(c)+B(c)f'(r)e$, then
$-h/r^2+|h|^2\cdot(-2r^{-3})e$, whose sum is the stated $p_1$; the $y'$
computation is identical with $e$ replaced by $-e$, giving $p_2$.

\emph{(5a) Bounds on the data.} By \eqref{eq:tpf-f}, \eqref{eq:tpf-B},
\eqref{eq:tpf-h}, $f(r)\le r\le3/2$ and $|q|\le C$,
\begin{equation}\label{eq:tpf-bounds}
  0<a\le B_{\mathrm{out}},\qquad
  \Bigl|\tfrac12(p_1+p_2)\Bigr|=\Bigl|q+\tfrac12f(r)DB(c)-\frac h{r^2}\Bigr|
   \le C+\tfrac34K_1+2C=:K_3,
\end{equation}
where $a\le B(c)f'(r)\le B_{\mathrm{out}}$ because $2|h|^2/r^3\ge0$ and
$f'\le1$; the positivity of $a$ is part of (5b).

\emph{(5b) $a\ge\delta/4$ for small $r$.} Assume in addition
\begin{equation}\label{eq:tpf-r2}
  \tfrac3{16}\sqrt r\le\tfrac14\qquad\text{and}\qquad 8C^2r\le\frac\delta8 .
\end{equation}
Then $f'(r)\ge3/4$, and $2|h|^2/r^3\le2(2Cr^2)^2/r^3=8C^2r\le\delta/8$ by
\eqref{eq:tpf-h}, so
\[
  a\ge\frac\delta2\cdot\frac34-\frac\delta8=\frac{3\delta}8-\frac\delta8
   =\frac\delta4>0 .
\]

\emph{(5c) The matrix inequality.} Let $\eta>0$. Apply \noderef{SumRuleJets}
with $N=3$, the open sets $O_1=O_2=O_3=B_1$, the functions $g_1,g_2,g_3$, the
test $\Psi$, the point $(x,y,z)$ and this $\eta$: all its hypotheses have just
been verified. It supplies symmetric continuous bilinear forms $X_1,X_2,X_3$ on
$E$ with, $\overline J^{2,+}$ denoting throughout the closed superjet of
\noderef{ClosedSuperjet}, as in the statements of both \noderef{SumRuleJets} and
\noderef{InfSubsolutionClosedSuperjet},
\begin{equation}\label{eq:tpf-jets}
  (p_1,X_1)\in\overline J^{2,+}U_i(x),\quad
  (p_2,X_2)\in\overline J^{2,+}U_i(y),\quad
  (-2p_0,X_3)\in\overline J^{2,+}(-2U_i)(z),
\end{equation}
and, writing $A=D^2\Psi(x,y,z)$ and $\bar A$ for the associated self-adjoint
operator,
\begin{equation}\label{eq:tpf-matrix}
  X_1[\xi_1]+X_2[\xi_2]+X_3[\xi_3]\le A[\xi]+\eta|\bar A\xi|^2
  \qquad\bigl(\xi=(\xi_1,\xi_2,\xi_3)\bigr).
\end{equation}
By Step 0 the function $U_i$ is an infinity subsolution on $B_1$ and $-2U_i$ is
an infinity subsolution on $B_1$. Therefore
\noderef{InfSubsolutionClosedSuperjet}, applied three times --- to $U_i$ at
$x\in B_1$ with $(p_1,X_1)$, to $U_i$ at $y\in B_1$ with $(p_2,X_2)$, and to
$-2U_i$ at $z\in B_1$ with $(-2p_0,X_3)$ --- gives
\[
  X_1[p_1]\ge0,\qquad X_2[p_2]\ge0,\qquad X_3[-2p_0]\ge0 .
\]
The last reads $4X_3[p_0]\ge0$ by bilinearity, so $X_3[p_0]\ge0$. Taking
$\xi=\sigma:=(p_1,p_2,p_0)$ in \eqref{eq:tpf-matrix},
\[
  0\le X_1[p_1]+X_2[p_2]+X_3[p_0]\le A[\sigma]+\eta|\bar A\sigma|^2 .
\]
Neither $A$ nor $\sigma$ depends on $\eta$, so letting $\eta\downarrow0$ yields
\begin{equation}\label{eq:tpf-Apos}
  A[\sigma]\ge0 .
\end{equation}

\emph{Step 6: the second variation is very negative.}
Because $\Psi$ is $C^2$ near $(x,y,z)$,
\[
  A[\sigma]=\frac{d^2}{d\tau^2}
   \Psi\bigl(x+\tau p_1,\ y+\tau p_2,\ z+\tau p_0\bigr)\Big|_{\tau=0}.
\]
Write $c(\tau),h(\tau),r(\tau)$ for the quantities induced by that line. Since
$p_1-p_2=2ae$ by \eqref{eq:tpf-grads},
\[
  (x+\tau p_1)-(y+\tau p_2)=re+2a\tau e=(r+2a\tau)e,
\]
so $r(\tau)=r+2a\tau$ for $|\tau|$ small (here $a>0$ and $r>0$), giving
$\dot r=2a$ and $\ddot r=0$. Also $c(\tau)=c+\tau\dot c$ with
$\dot c=\tfrac12(p_1+p_2)$, and $h(\tau)=h+\tau\dot h$ with
\[
  \dot h=p_0-\dot c
   =\Bigl(q-\frac h{r^2}\Bigr)-\Bigl(q+\tfrac12f(r)DB(c)-\frac h{r^2}\Bigr)
   =-\tfrac12f(r)DB(c),
  \qquad|\dot h|\le\tfrac12rK_1 ,
\]
using $f(r)\le r$ and $c\in\overline B(0,3/4)$. The affine part
$\ell_i(x+\tau p_1)+\ell_i(y+\tau p_2)-2\ell_i(z+\tau p_0)$ is affine in $\tau$
and contributes $0$.

\emph{(6a) The penalty term.} With $u(\tau)=|h(\tau)|^2$ and
$v(\tau)=r(\tau)^{-2}$ we have $u(0)=|h|^2$, $\dot u(0)=2\langle h,\dot
h\rangle$, $\ddot u(0)=2|\dot h|^2$, $v(0)=r^{-2}$, $\dot v(0)=-2\dot r r^{-3}$,
$\ddot v(0)=6\dot r^2r^{-4}$, so
\[
  \frac{d^2}{d\tau^2}\frac{|h(\tau)|^2}{r(\tau)^2}\Big|_{\tau=0}
   =\frac{2|\dot h|^2}{r^2}-\frac{8\dot r\langle h,\dot h\rangle}{r^3}
     +\frac{6\dot r^{\,2}|h|^2}{r^4}.
\]
By $|\dot h|\le rK_1/2$, $|h|\le2Cr^2$, $\dot r=2a\le2B_{\mathrm{out}}$, the
three terms are bounded in absolute value by $K_1^2/2$, $16B_{\mathrm{out}}CK_1$
and $96B_{\mathrm{out}}^2C^2$ respectively.

\emph{(6b) The localizing term.} Since $c(\tau)$ and $r(\tau)$ are affine in
$\tau$,
\[
  \frac{d^2}{d\tau^2}\bigl(B(c(\tau))f(r(\tau))\bigr)\Big|_{\tau=0}
   =f(r)D^2B(c)[\dot c]+2f'(r)\dot r\,\langle DB(c),\dot c\rangle
     +B(c)f''(r)\dot r^{\,2}.
\]
The first two terms are bounded in absolute value by $\tfrac32K_2K_3^2$ and
$4B_{\mathrm{out}}K_1K_3$ by \eqref{eq:tpf-f} and \eqref{eq:tpf-bounds}. The
third is
\[
  B(c)\cdot\Bigl(-\frac3{32\sqrt r}\Bigr)\cdot4a^2
   =-\frac{3B(c)a^2}{8\sqrt r}\le-\frac{3}{8\sqrt r}\cdot\frac\delta2
      \cdot\frac{\delta^2}{16}=-\frac{3\delta^3}{256\sqrt r},
\]
using $B(c)\ge\delta/2$ and $a\ge\delta/4$ from (5b).

Collecting, with
\[
  M:=\frac{K_1^2}2+16B_{\mathrm{out}}CK_1+96B_{\mathrm{out}}^2C^2
    +\tfrac32K_2K_3^2+4B_{\mathrm{out}}K_1K_3 ,
\]
a constant depending only on $C$, $\delta$, $K_1$, $K_2$, $K_3$ and
$B_{\mathrm{out}}$ --- and so on nothing that varies with $k$ --- we obtain
\begin{equation}\label{eq:tpf-neg}
  A[\sigma]\le M-\frac{3\delta^3}{256\sqrt r} .
\end{equation}

\emph{Step 7: the contradiction.}
Let
\[
  r_0=\min\Bigl\{\frac14,\ \frac1{16C},\ \frac{16}{9},\
    \frac{\delta}{64C^2},\
    \Bigl(\frac{3\delta^3}{256(M+1)}\Bigr)^{2}\Bigr\}>0 ,
\]
a positive number independent of $k$. If $r<r_0$ then: $r<1/2$ and
$2Cr^2<2C\cdot\frac1{16C}\cdot r<\frac14$, so \eqref{eq:tpf-r1} holds;
$\tfrac3{16}\sqrt r\le\tfrac3{16}\cdot\tfrac43=\tfrac14$ and
$8C^2r\le8C^2\delta/(64C^2)=\delta/8$, so \eqref{eq:tpf-r2} holds; and
$\sqrt r<3\delta^3/(256(M+1))$, so
$3\delta^3/(256\sqrt r)>M+1$ and \eqref{eq:tpf-neg} gives $A[\sigma]<-1<0$,
contradicting \eqref{eq:tpf-Apos}.

So it remains to produce one index $k$ whose maximizing triple has $r<r_0$. Put
$\varepsilon_0=\delta\nu r_0/8>0$ and choose $k$ with
$i_k\ge\max\bigl(\mathcal E(1),\mathcal E(\varepsilon_0)\bigr)$, which is
possible because $i_k\to\infty$. For this $k$ all the hypotheses used in Steps
3--6 hold, and \eqref{eq:tpf-rsmall} with $\varepsilon=\varepsilon_0$ gives
$r<8\varepsilon_0/(\delta\nu)=r_0$. This is the required contradiction, so the
conclusion cannot fail.
\end{proof}
